\chapter{Output Reductions: From Fields to Physical Quantities}
\label{ch:reductions}

\chapterorient{Every analysis in this book ends the same way. A driver has
worked hard to produce a \emph{family of fields}---one potential per terminal,
one scattered field per frequency, one eigenvector per resonant mode---and now
that family must be boiled down to the single number, or small table, the
engineer actually asked for: a capacitance, an S-matrix, a list of resonant
frequencies and their quality factors. This last stage is the \texttt{reduce}
of the skeleton (\cref{ch:situating}), and it is where the physics finally
re-enters: a Gram matrix \emph{is} a capacitance, a projection \emph{is} an
S-parameter, an eigenvalue's imaginary part \emph{is} a damping rate. The
surprise of this chapter---the same surprise the whole book keeps finding---is
that the dozen physical outputs Palace can produce are built from only
\emph{three} reduction shapes. We meet all three in detail, compute two of them
by hand, and close Part~IV by handing the assembled machinery forward to
Part~V, where the modalities compose it.}

\section{The setting: a solution family, and the answer hidden in it}
\label{sec:red-setting}

By the time we reach the reduce stage, the hard work is done. \Cref{ch:situating}
drew the skeleton---\texttt{config} $\to$ \texttt{assemble} $\to$ \texttt{solve}
$\to$ \texttt{reduce} $\to$ \texttt{product}---and Parts~II through~IV have now
filled in every box but the last. Assembly turns the continuous physics into
operators (\cref{ch:assembly}); the solve stage, in one of its four corners
(\cref{ch:situating}), turns those operators into fields. What the solve hands
to the reduce stage is never a single field. It is a \emph{family}:

\begin{itemize}[nosep]
  \item the electrostatic and magnetostatic drivers produce one solved field
  \emph{per excitation}---per terminal held at unit voltage, per source carrying
  unit current (\cref{ch:electrostatics}, \cref{ch:magnetostatics});
  \item the driven driver produces one scattered field \emph{per frequency} in
  the sweep, and within each, one column \emph{per driven port}
  (\cref{ch:driven});
  \item the eigenmode and boundary-mode drivers produce one eigenvector
  \emph{per converged mode} (\cref{ch:eigenmodes}, \cref{ch:waveports}).
\end{itemize}

The reduce stage takes that family and collapses it to the product. The
physical answer is, in a precise sense, already \emph{present} in the family---it
is a pairing, a projection, or a per-item readout away. What the reduction does
is extract it (\cref{fig:red-pipeline}).

\begin{figure}[t]
\centering
\begin{tikzpicture}[node distance=8mm and 12mm]
  \node[stage] (solve) {\textbf{solve}};
  \node[data, right=14mm of solve, align=center, text width=26mm] (fam)
    {solution family\\[1pt]\footnotesize $\{\vx_i\}$, $\{\vE_j(\omega)\}$,
     $\{(\lambda_m,\vx_m)\}$};
  \node[stage, right=14mm of fam, fill=simBgGold, draw=simGold] (red)
    {\textbf{reduce}};
  \node[product, right=12mm of red, align=center, text width=24mm] (prod)
    {physical\\product\\[1pt]\footnotesize $C$, $S$, $(f,Q)$, $\dots$};
  \draw[flow] (solve) -- (fam);
  \draw[flow] (fam) -- node[above,font=\scriptsize]{map-then-collect} (red);
  \draw[flow] (red) -- (prod);
  \node[font=\scriptsize\itshape, simGray, below=2mm of red, align=center]
    {embarrassingly parallel:\\no inter-element state};
\end{tikzpicture}
\caption[The family-to-product pipeline]{The reduce stage in one picture. The
solve produces not a single field but a \emph{family}---one member per
excitation, per frequency, or per mode. The reduction maps a per-element
computation over that family and collects the results into the small physical
product. Because no element depends on another, every reduction in this chapter
is a \texttt{map}-then-collect (or \texttt{map}-then-reduce), the embarrassingly
parallel combinator of \cref{ch:combinators}: there is no state threaded between
elements, so the work is independent and reorderable.}
\label{fig:red-pipeline}
\end{figure}

That structural observation---\emph{every} reduction is a map over the family
with no inter-element state---is worth stating up front, because it is what
makes the reductions so uniform. None of them is a \texttt{fold}; none threads a
running accumulator from one family member to the next. Each is the
embarrassingly parallel \texttt{map} combinator of \cref{ch:combinators},
followed by a collect into a matrix or a table. In the framework's vocabulary
they are \emph{pure value-producing reductions}: no \code{Solve} monad, no
convergence predicate, no carried state---post-processing readouts that take a
solved family and return a value.

\begin{keyidea}
The reduce stage turns a \emph{solution family} into the user-facing physical
quantity. Every reduction is a \texttt{map} over the family with \emph{no
inter-element state}---embarrassingly parallel, never a fold
(\cref{ch:combinators}). What varies between reductions is only the per-element
computation and how the results are collected. There are exactly three
collection \emph{shapes}, and the rest of the chapter is those three.
\end{keyidea}

\Cref{ch:situating} named the three shapes when it first sketched the framework;
here we give each its full treatment. They are, in order: the \emph{rank-2
symmetric Gram} (capacitance, inductance); the \emph{rank-2 port-projection}
(S-parameters); and the \emph{rank-1 scalar-or-mode table} (energy and
participation, eigenfrequency and $Q$, waveguide modes). \Cref{fig:red-taxonomy}
is the hero figure of the chapter: the three shapes side by side, the taxonomy
that organizes everything that follows.

\begin{figure}[p]
\centering
\begin{tikzpicture}[font=\footnotesize]
  % ===== (a) symmetric Gram =====
  \begin{scope}[shift={(0,0)}]
    \foreach \i in {0,1,2}{
      \foreach \j in {0,1,2}{
        \node[draw=simGray, semithick,
          fill={\ifnum\i=\j simBgRust\else simBgBlue\fi},
          minimum size=7mm, inner sep=0pt]
          (g\i\j) at (\j*0.72, -\i*0.72) {};}}
    % mirror diagonal
    \draw[densely dashed, simRust, semithick] (-0.5,0.5) -- (1.94,-1.94);
    % an off-diagonal pair, mirrored
    \node[font=\scriptsize, simBlue] at (1.44,0.0) {$G_{02}$};
    \node[font=\scriptsize, simBlue] at (0.0,-1.44) {$G_{20}$};
    \node[font=\scriptsize\bfseries, align=center] at (0.72,-2.55)
      {\textsc{(a)} rank-2 Gram};
    \node[font=\scriptsize, align=center] at (0.72,-3.05)
      {$G_{ij}=w(i,j)\,\vx_j^{\!\top}\Kop\,\vx_i$};
    \node[font=\scriptsize, align=center] at (0.72,-3.5)
      {\emph{symmetric}: $G_{ij}=G_{ji}$};
    \node[font=\scriptsize\itshape, text=simRust, align=center] at (0.72,-4.0)
      {capacitance,\\inductance};
  \end{scope}

  % ===== (b) non-symmetric projection =====
  \begin{scope}[shift={(4.6,0)}]
    \foreach \i in {0,1,2}{
      \foreach \j in {0,1,2}{
        \node[draw=simGray, semithick, fill=simBgGreen,
          minimum size=7mm, inner sep=0pt]
          (p\i\j) at (\j*0.72, -\i*0.72) {};}}
    \node[font=\scriptsize, text=simGreen] at (0.72,0.6) {drive port $j\to$};
    \node[font=\scriptsize, rotate=90, text=simGreen] at (-0.72,-0.72)
      {receiver port $i\to$};
    \node[font=\scriptsize\bfseries, align=center] at (0.72,-2.55)
      {\textsc{(b)} rank-2 projection};
    \node[font=\scriptsize, align=center] at (0.72,-3.05)
      {$S_{ij}=\sigma_{ij}\,(\inner{\vs_i}{\vE_j}-\delta_{ij})$};
    \node[font=\scriptsize, align=center] at (0.72,-3.5)
      {\emph{not} symmetric, complex};
    \node[font=\scriptsize\itshape, text=simGreen, align=center] at (0.72,-4.0)
      {S-parameters};
  \end{scope}

  % ===== (c) rank-1 table =====
  \begin{scope}[shift={(9.2,0)}]
    \foreach \i in {0,1,2,3}{
      \node[draw=simGray, semithick, fill=simBgGold,
        minimum width=22mm, minimum height=5.5mm, inner sep=0pt]
        (r\i) at (0, -\i*0.66) {};}
    \node[font=\scriptsize] at (0,0)     {mode 1: $f_1,\,Q_1$};
    \node[font=\scriptsize] at (0,-0.66) {mode 2: $f_2,\,Q_2$};
    \node[font=\scriptsize] at (0,-1.32) {mode 3: $f_3,\,Q_3$};
    \node[font=\scriptsize] at (0,-1.98) {$\vdots$};
    \node[font=\scriptsize\bfseries, align=center] at (0,-2.55)
      {\textsc{(c)} rank-1 table};
    \node[font=\scriptsize, align=center] at (0,-3.05)
      {one row per item};
    \node[font=\scriptsize, align=center] at (0,-3.5)
      {a \texttt{map}, no pairing};
    \node[font=\scriptsize\itshape, text=simGold, align=center] at (0,-4.0)
      {energy$+p$, $f{+}Q$,\\waveguide modes};
  \end{scope}
\end{tikzpicture}
\caption[The three reduction shapes]{The three reduction shapes---the chapter's
organizing taxonomy, first sketched in \cref{ch:situating}. \textsc{(a)} A
\emph{rank-2 symmetric Gram} pairs a solution family with itself through an
operator, $G_{ij}=w(i,j)\,\vx_j^{\!\top}\Kop\,\vx_i$; its symmetry (dashed
mirror) is structural, so only the upper triangle is computed and the lower is
copied ($G_{02}=G_{20}$). \textsc{(b)} A \emph{rank-2 projection} also yields a
matrix, but it projects each driven field onto a basis of port modes; its two
indices play different roles, it is complex, and it is in general \emph{not}
symmetric. \textsc{(c)} A \emph{rank-1 table} is a one-dimensional list, one row
per item, with no pairing at all---a pure \texttt{map}. The single most
over-unified point in the design is the line between \textsc{(a)} and
\textsc{(b)}: two matrix-shaped results are \emph{not} the same reduction
(\cref{pit:gram-not-proj}).}
\label{fig:red-taxonomy}
\end{figure}

\section{Shape 1: the symmetric Gram reduction}
\label{sec:gram}

The first shape is a \emph{Gram matrix}: a symmetric matrix of operator-weighted
inner products formed by pairing a solution family with itself. It is the
reduction behind both capacitance and inductance, and it is the cleanest
instance in the book of ``one verb, two products.''

\subsection{The verb}

Given a solution family $\vx_0,\dots,\vx_{m-1}$, the assembled \emph{energy
operator} $\Kop$ (symmetric positive-definite), and a scalar \emph{weight}
function $w(i,j)$, the Gram reduction computes the $m\times m$ matrix
\begin{equation}
  \boxed{\;G_{ij} \;=\; w(i,j)\;\vx_j^{\!\top}\,\Kop\,\vx_i\;}
  \label{eq:gram-entry}
\end{equation}
The entry $G_{ij}$ is an \emph{operator-weighted bilinear form} of family
members $i$ and $j$, scaled by the weight. In the framework this verb is
\code{gram\_reduce}, with the signature and body
\begin{lstlisting}[style=l4style]
gram_reduce :: LinOp -> [Tensor] -> (Int -> Int -> Scalar) -> Matrix
gram_reduce k xs w =
  symmetric_from_upper                          -- mirror lower triangle from upper
    [ [ w i j * entry k xs i j | j <- [i .. m-1] ]   -- map over upper-triangle pairs
      | i <- [0 .. m-1] ]
  where
    m = length xs
    entry k xs i j
      | i == j    = matrix_weighted_norm (xs!!i) k   -- diagonal: x_i^T K x_i
      | otherwise = bilinear_form (xs!!j) k (xs!!i)  -- off-diag:  x_j^T K x_i
\end{lstlisting}
Two structural facts are baked into this body. First, each grid entry is
\emph{independent}: the upper-triangle comprehension is a \texttt{map} over
family pairs, with no state threaded between entries---the embarrassingly
parallel pattern of \cref{sec:red-setting}. Second, the matrix is \emph{symmetric
by construction}: the body computes only the upper triangle and mirrors the lower
from it (\code{symmetric\_from\_upper}). The symmetry is not checked after the
fact; it is guaranteed because $\Kop$ is symmetric, so $\vx_j^{\!\top}\Kop\vx_i
=\vx_i^{\!\top}\Kop\vx_j$, and the verb exploits it to halve the work.

\begin{notationbox}
The diagonal and off-diagonal entries call two slightly different lower-layer
verbs. On the diagonal, $G_{ii}=w(i,i)\,\vx_i^{\!\top}\Kop\vx_i$ is a
\emph{self}-pairing---an operator-weighted squared norm, $\code{matrix\_weighted\_norm}$
(its radicand). Off the diagonal, $G_{ij}=w(i,j)\,\vx_j^{\!\top}\Kop\vx_i$ is a
\emph{cross}-pairing of two different members, the bilinear form
$\code{bilinear\_form}$. They are the same expression $\vx_j^{\!\top}\Kop\vx_i$
specialized to $i=j$ and $i\neq j$; the verb names them separately only because
the self-case has a cheaper, symmetric implementation.
\end{notationbox}

\subsection{Two products, one weight apart}

Here is the payoff. \emph{Capacitance and inductance are the same reduction.}
They call the identical verb \code{gram\_reduce} on the identical operator-weighted
bilinear form; they differ in exactly one place---the scalar weight $w(i,j)$---and
in what the family members physically are.

\paragraph{Capacitance.} The electrostatic driver holds each conductor $i$ at
unit voltage and solves for the potential field $\vx_i=V_i$
(\cref{ch:electrostatics}). The operator $\Kop$ is the
$\varepsilon$-weighted energy operator (the discrete form of
$\tfrac12\varepsilon\abs{\vE}^2$, \cref{ch:maxwell}, eq.~\eqref{eq:energy-density}).
The weight is \emph{unity}:
\begin{equation}
  w(i,j) = 1,
  \qquad
  C_{ij} = V_j^{\!\top}\,\Kop\,V_i .
  \label{eq:cap-gram}
\end{equation}
The diagonal entry $C_{ii}=V_i^{\!\top}\Kop V_i$ is exactly twice the electric
field energy stored when conductor $i$ alone is at unit voltage---the energy form
$C_{ii}=2U_e/V^2$ with $V=1$ that \cref{ch:maxwell} promised
(eq.~\eqref{eq:energy-density}: $C=2U/V^2$). The off-diagonal $C_{ij}$ is the
\emph{mutual} capacitance, the charge induced on conductor $i$ by a unit voltage
on conductor $j$.

\paragraph{Inductance.} The magnetostatic driver drives each source $i$ with
unit current and solves for the vector potential field $\vx_i=\vA_i$
(\cref{ch:magnetostatics}). The operator $\Kop$ is now the $\nu$-weighted
curl--curl energy operator (the magnetic energy $\tfrac{1}{2\mu}\abs{\vB}^2$).
The verb, the bilinear form, the symmetry, the map---all identical. The only
change is the weight, which now \emph{normalizes by the currents}:
\begin{equation}
  w(i,j) = \frac{1}{I_i\,I_j},
  \qquad
  L_{ij} = \frac{1}{I_i I_j}\,\vA_j^{\!\top}\,\Kop\,\vA_i .
  \label{eq:ind-gram}
\end{equation}
The diagonal $L_{ii}$ is the self-inductance ($L=2U_m/I^2$, \cref{ch:maxwell});
the off-diagonal $L_{ij}$ is the mutual inductance between loops $i$ and $j$.

Why does inductance need the current normalization that capacitance does not? An
honest one-line answer: in electrostatics the excitation is a \emph{voltage} and
we choose it to be unity, so dividing by it does nothing; in magnetostatics the
natural excitation is a \emph{current}, and the energy bilinear form scales with
the square of the driving current, so to recover the geometry-only quantity
$L=2U_m/I^2$ we must divide out $I_iI_j$. The weight is exactly the
``per-unit-excitation'' rescaling, trivial in one case and nontrivial in the
other. \Cref{fig:gram-energy} draws the energy picture that underlies the
diagonal.

\begin{figure}[t]
\centering
\begin{tikzpicture}[font=\footnotesize, node distance=7mm]
  % the field-energy picture for the Gram diagonal
  \node[data, align=center, text width=22mm] (vi) {field $\vx_i$\\\footnotesize
    (unit excitation $i$)};
  \node[opbox, minimum width=20mm, right=10mm of vi] (Kx) {$\Kop\,\vx_i$};
  \node[opbox, fill=simBgRust, draw=simRust, minimum width=30mm, right=10mm of Kx]
    (form) {$\vx_i^{\!\top}\Kop\,\vx_i = 2U$};
  \node[product, align=center, text width=26mm, right=10mm of form] (out)
    {$G_{ii}=w(i,i)\cdot 2U$\\[1pt]\footnotesize $=C_{ii}$ or $L_{ii}$};
  \draw[flow] (vi) -- node[above,font=\scriptsize]{apply $\Kop$} (Kx);
  \draw[flow] (Kx) -- node[above,font=\scriptsize]{$\inner{\cdot}{\vx_i}$} (form);
  \draw[flow] (form) -- node[above,font=\scriptsize]{$\times\,w$} (out);
  \node[font=\scriptsize\itshape, simGray, below=2mm of form, align=center]
    {$U=\tfrac12\!\int\!\varepsilon\abs{\vE}^2$ (elec.)\\
     or $\tfrac{1}{2\mu}\!\int\!\abs{\vB}^2$ (mag.)};
  \node[font=\scriptsize\itshape, simRust, below=10mm of out, align=center]
    {$w=1$: capacitance\\$w=1/I_i^2$: inductance};
\end{tikzpicture}
\caption[The Gram diagonal is field energy]{The energy picture behind the Gram
diagonal. Applying the energy operator $\Kop$ to a solved field $\vx_i$ and
pairing the result back with $\vx_i$ yields $\vx_i^{\!\top}\Kop\vx_i$, which is
\emph{twice the stored field energy} $2U$ (\cref{ch:maxwell},
eq.~\eqref{eq:energy-density}). Scaling by the weight $w$ gives the physical
diagonal: with $w=1$ this is the self-capacitance $C_{ii}=2U_e/V^2$ at unit
voltage; with $w=1/I_i^2$ it is the self-inductance $L_{ii}=2U_m/I^2$ at unit
current. The off-diagonal repeats the picture with two \emph{different} fields,
$\vx_j^{\!\top}\Kop\vx_i$, giving mutual capacitance or inductance.}
\label{fig:gram-energy}
\end{figure}

\begin{keyidea}
The \textbf{symmetric Gram reduction} $G_{ij}=w(i,j)\,\vx_j^{\!\top}\Kop\vx_i$ is
one verb (\code{gram\_reduce}) that produces \emph{two} physical products.
Capacitance uses unit weight $w=1$ on the per-terminal potential family;
inductance uses the current-normalized weight $w=1/(I_iI_j)$ on the per-source
vector-potential family. Same operator-weighted bilinear form, same structural
symmetry ($G_{ij}=G_{ji}$, computed on the upper triangle and mirrored), same
embarrassingly parallel map---differing only in a scalar weight and in what the
family physically is. The diagonal is twice the stored field energy; the
off-diagonal is the mutual coupling.
\end{keyidea}

\subsection{Worked example: a two-terminal capacitance}

\begin{workedexample}{A $2\times 2$ capacitance from a symmetric Gram}{cap2x2}
A two-conductor electrostatic problem has been solved. The driver held conductor
$1$ at unit voltage (conductor $2$ grounded) and solved for the potential field
$V_1$; then it held conductor $2$ at unit voltage and solved for $V_2$. The
discrete fields and the assembled energy operator $\Kop$ are given; we are to
build the $2\times 2$ capacitance matrix.

Suppose evaluating the operator-weighted bilinear form $\vx_j^{\!\top}\Kop\vx_i$
on the three needed pairs gives the raw numbers (in femtofarads, with unit
voltage so $w\equiv 1$):
\[
  V_1^{\!\top}\Kop V_1 = 4.20,\qquad
  V_2^{\!\top}\Kop V_2 = 3.60,\qquad
  V_2^{\!\top}\Kop V_1 = -1.10 .
\]
We compute the upper triangle and \emph{mirror} the off-diagonal---we never
recompute $V_1^{\!\top}\Kop V_2$, because the symmetry of $\Kop$ guarantees it
equals $V_2^{\!\top}\Kop V_1$:
\[
  C_{11}=4.20,\qquad
  C_{22}=3.60,\qquad
  C_{12}=C_{21}=-1.10 .
\]
The capacitance matrix is therefore
\[
  \mat{C} \;=\;
  \begin{pmatrix} 4.20 & -1.10 \\ -1.10 & 3.60 \end{pmatrix}\ \si{\femto\farad}.
\]

\emph{Interpretation.} The diagonal entries are the \emph{self}-capacitances:
held alone at unit voltage, conductor $1$ stores $\tfrac12 C_{11}V^2$ of field
energy, conductor $2$ slightly less. The off-diagonal is the \emph{mutual}
capacitance, and its \emph{sign is negative}---a universal feature of the
Maxwell capacitance matrix. Raising conductor $2$'s voltage draws charge of the
\emph{opposite} sign onto conductor $1$ through their mutual field, so the
induced charge $Q_1=C_{11}V_1+C_{12}V_2$ falls as $V_2$ rises; the off-diagonal
coupling is an inducement of opposite charge, hence $C_{12}<0$. The matrix is
symmetric ($C_{12}=C_{21}$) because the underlying energy operator is symmetric:
the charge that conductor $2$ induces on conductor $1$ per volt equals the charge
conductor $1$ induces on conductor $2$ per volt---a reciprocity that the Gram
construction makes automatic, not coincidental.

\emph{The inductance contrast.} Had this been a magnetostatic problem with the
same raw bilinear values but driving currents $I_1=2$, $I_2=5$ (in amperes), the
weight $w(i,j)=1/(I_iI_j)$ would have rescaled each entry:
$L_{11}=4.20/4=1.05$, $L_{22}=3.60/25=0.144$, and
$L_{12}=L_{21}=-1.10/(2\cdot 5)=-0.110$ (in the corresponding units). Same verb,
same symmetry, same fields-in; one weight closure swapped, and the energy form
becomes an inductance instead of a capacitance.
\end{workedexample}

\section{Shape 2: the port-projection reduction}
\label{sec:proj}

The second shape is also a matrix indexed by two families---and it is tempting to
file it under the first. Resist. The S-parameter matrix is a \emph{projection},
not a Gram, and the difference is not cosmetic.

\subsection{The verb}

The driven driver solves the time-harmonic system at each frequency $\omega$,
once per driven port, producing a family of complex scattered fields
$\vE_j(\omega)$---one column per drive port $j$ (\cref{ch:driven}). The
S-parameter reduction projects each of those fields onto every \emph{receiver}
port's mode $\vs_i$, subtracts a self-reflection on the diagonal, and scales by
an impedance/de-embedding factor:
\begin{equation}
  \boxed{\;S_{ij} \;=\; \sigma_{ij}\,\bigl(\inner{\vs_i}{\vE_j(\omega)} - \delta_{ij}\bigr)\;}
  \label{eq:sparam-entry}
\end{equation}
where $\inner{\vs_i}{\vE_j}$ is the linear projection of the drive-$j$ field onto
receiver mode $i$, $\delta_{ij}$ is the Kronecker delta (the self-term, $1$ when
$i=j$ and $0$ otherwise), and $\sigma_{ij}$ is the port-kind scale (an impedance
ratio for lumped ports, a de-embedding phase for wave ports). The framework verb
is \code{sparameter\_reduce}:
\begin{lstlisting}[style=l4style]
sparameter_reduce :: [PortMode] -> [(Int, Tensor)] -> Matrix
sparameter_reduce ports family =
  matrix_from_columns
    [ [ scale ports i j * (project (ports!!i) e - selfterm i j)  -- entry S_ij
        | i <- [0 .. p-1] ]                                      -- over receivers i
      | (j, e) <- family ]                                       -- one column per drive port j
  where
    p             = length ports
    project s e   = port_dot s e            -- linear functional s_i . E
    selfterm i j  = if i == j then 1 else 0 -- the -1 self-reflection on the diagonal
    scale i j     = port_scale (ports!!i) (ports!!j)  -- impedance / de-embed
\end{lstlisting}
The reduction is applied once per swept frequency; the $\omega$ index rides as an
outer family axis that the driven composition root owns, so the verb itself is
the per-port reduction at a single $\omega$.

\subsection{Why this is rank-2 like a Gram but deliberately not a Gram}

Both \code{gram\_reduce} and \code{sparameter\_reduce} produce square matrices
indexed by a port-like family. That shared \emph{output rank} is the whole of
their resemblance, and reading more into it is the single most over-unified
mistake the design could make.

\begin{pitfall}
\label{pit:gram-not-proj}
\textbf{Two reductions producing matrices are not the same reduction.} It is
sorely tempting to unify \code{gram\_reduce} and \code{sparameter\_reduce} under
one ``matrix reduction'' verb, since both yield a $p\times p$ array indexed by
ports. Three structural facts forbid it.
\begin{itemize}[nosep]
  \item \emph{The Gram is a self-pairing; the projection is not.} The Gram pairs
  a family \emph{with itself} through an operator, $\vx_j^{\!\top}\Kop\vx_i$. The
  S-matrix projects one family (the solved drive fields $\vE_j$) onto a
  \emph{different} basis (the port modes $\vs_i$)---two families in different
  roles, not one family paired with itself.
  \item \emph{The Gram is symmetric; the projection is not.} The Gram's symmetry
  $G_{ij}=G_{ji}$ is structural (\cref{sec:gram}) and the verb \emph{computes
  only the upper triangle}. The S-matrix has no such guarantee---every entry is
  computed independently---and its near-symmetry $S_{ij}\approx S_{ji}$, where it
  holds at all, is reciprocity \emph{physics}, not a construction the reduction
  imposes. There is no \code{symmetric\_from\_upper} in the projection body.
  \item \emph{The Gram is real and energy-valued; the projection is complex and
  has an inhomogeneous diagonal.} S-parameters are complex (magnitude and phase),
  and the diagonal carries the $-\delta_{ij}$ self-reflection term, which has no
  analogue in the Gram. Forcing one verb would either smuggle a false symmetry
  into the S-matrix or burden the Gram with a self-term and complex arithmetic it
  does not need.
\end{itemize}
The right move---which the framework makes---is two verbs that happen to share an
output rank. \emph{Sharing a shape is not sharing a meaning.}
\end{pitfall}

The port modes $\vs_i$ themselves are not produced by the driven driver; they
come from the \emph{boundary-mode} analysis (\cref{ch:waveports}), which solves a
small two-dimensional eigenproblem on each port's cross-section to find its
propagation modes. The driven reduction consumes those modes as the basis it
projects onto. \Cref{fig:proj} draws the projection.

\begin{figure}[t]
\centering
\begin{tikzpicture}[font=\footnotesize]
  % the solved field as an arrow in "field space"
  \draw[->, simGray] (0,0) -- (0,3.2) node[above,font=\scriptsize]{};
  \draw[->, simGray] (0,0) -- (4.6,0) node[right,font=\scriptsize]{};
  % port mode basis vectors
  \draw[-{Stealth[length=2.4mm]}, very thick, simGreen] (0,0) -- (3.4,0.0)
    node[right,font=\scriptsize,simGreen]{$\vs_1$ (port 1 mode)};
  \draw[-{Stealth[length=2.4mm]}, very thick, simGreen] (0,0) -- (0.0,2.6)
    node[above,font=\scriptsize,simGreen]{$\vs_2$ (port 2 mode)};
  % the solved driven field E_j
  \draw[-{Stealth[length=2.6mm]}, very thick, simRust] (0,0) -- (2.7,1.7)
    node[above right,font=\scriptsize,simRust]{$\vE_j(\omega)$ (drive $j$)};
  % projections (dashed drops)
  \draw[densely dashed, simBlue] (2.7,1.7) -- (2.7,0)
    node[midway,right,font=\scriptsize,simBlue]{};
  \draw[densely dashed, simBlue] (2.7,1.7) -- (0,1.7);
  \node[font=\scriptsize, simBlue] at (2.7,-0.28) {$\inner{\vs_1}{\vE_j}$};
  \node[font=\scriptsize, simBlue] at (-0.6,1.7) {$\inner{\vs_2}{\vE_j}$};
  % annotation
  \node[font=\scriptsize\itshape, simGray, align=left] at (3.2,2.6)
    {each projection\\$=$ one $S_{\bullet j}$\\(before $-\delta$, scale)};
\end{tikzpicture}
\caption[Port projection of a driven field]{The S-parameter reduction as a
\emph{linear projection}. The driven driver's solved field $\vE_j(\omega)$ (rust)
for drive port $j$ is projected onto the basis of \emph{port modes} $\vs_i$
(green, supplied by the boundary-mode analysis, \cref{ch:waveports}). Each
projection $\inner{\vs_i}{\vE_j}$ (blue drop) becomes one column entry
$S_{ij}$---after subtracting the diagonal self-reflection $\delta_{ij}$ and
applying the impedance/de-embed scale $\sigma_{ij}$. Unlike the Gram of
\cref{fig:gram-energy}, the two index families are \emph{different} (drive fields
vs.\ port modes), the operation is a projection rather than an energy
self-pairing, and the result is complex and not symmetric.}
\label{fig:proj}
\end{figure}

\section{Shape 3: the rank-1 scalar and mode tables}
\label{sec:tables}

The third shape is the simplest in form and the most populous in membership: a
\emph{table}---a one-dimensional list with one row per item, each row a few
scalars (or, in one case, a few small fields). There is no matrix, no pairing of
a family with itself, just a \texttt{map} over a collection producing one tuple
per element. Three distinct products live here, and we take them in turn.

\subsection{Energy and participation, per domain}
\label{sec:energy-table}

The most driver-agnostic reduction in the book asks, of \emph{any} field-bearing
analysis: where does the field energy sit? Given a solved field and a configured
map of domain-restricted energy operators $\{i\mapsto \mat{M}_i\}$, the
\emph{domain-energy} reduction computes, per domain, the energy in that domain
and the fraction of the total it represents:
\begin{equation}
  \mathrm{energy}_i \;=\; \tfrac12\,\inner{\text{field}}{\mat{M}_i\,\text{field}},
  \qquad
  p_i \;=\; \frac{\mathrm{energy}_i}{e_{\text{total}}},
  \label{eq:energy-part}
\end{equation}
where $\mat{M}_i$ is the energy operator restricted to domain $i$ (an SPD form,
the domain-local piece of the global energy operator of \cref{ch:maxwell},
eq.~\eqref{eq:energy-density}) and $e_{\text{total}}$ is the whole-domain total.
The participation factor $p_i$ is the headline number---it says what fraction of
the electric (or magnetic) field energy lives in each material region, the
quantity a designer inspects to see whether energy is concentrating where it
should. The verb is \code{domain\_energy\_reduce}; it is a \texttt{map} over the
domain map with no inter-domain state.

This reduction is \emph{driver-agnostic}: it cares only that a field exists and
that domain operators are configured, so it runs identically as a post-process of
the electrostatic, eigenmode, driven, or transient analyses. It runs once per
field kind---electric ($\tfrac12\inner{\vE}{\mat{M}_i\vE}$) and magnetic
($\tfrac12\inner{\vB}{\mat{M}_i\vB}$)---producing one table per kind. The field
may be complex (in the driven and eigenmode cases), but the energy form is a real
nonnegative number, so the table is real. When the configured domains
\emph{partition} the field's support, the participations sum to one,
$\sum_i p_i = 1$; this is the natural sanity check on the table.
\Cref{fig:energy-bars} shows a typical participation table as a bar chart.

\begin{figure}[t]
\centering
\begin{tikzpicture}
\begin{axis}[
    width=0.82\linewidth, height=46mm,
    ybar, bar width=10mm,
    ymin=0, ymax=0.7,
    ylabel={participation $p_i$}, ylabel style={font=\scriptsize},
    xlabel={material domain}, xlabel style={font=\scriptsize},
    symbolic x coords={substrate, dielectric, air, metal},
    xtick=data, xticklabel style={font=\scriptsize},
    ytick={0,0.2,0.4,0.6}, yticklabel style={font=\scriptsize},
    nodes near coords, nodes near coords style={font=\scriptsize},
    every node near coord/.append style={/pgf/number format/fixed,
      /pgf/number format/precision=2},
    enlarge x limits=0.18,
    axis lines=left,
  ]
  \addplot[fill=simBgGold, draw=simGold] coordinates {
    (substrate,0.52) (dielectric,0.31) (air,0.13) (metal,0.04)};
\end{axis}
\end{tikzpicture}
\caption[Energy participation per domain]{An electric-energy participation table,
drawn as a bar chart. Each bar is one row of the rank-1 table: the fraction
$p_i=\mathrm{energy}_i/e_{\text{total}}$ of the total field energy stored in
material domain $i$ (\cref{eq:energy-part}). Here most of the energy sits in the
substrate and dielectric, very little in the surrounding air, and almost none in
the metal---a typical signature for a planar device. The four participations sum
to $1.00$ because the configured domains partition the field's support; that sum
is the table's natural sanity check. The reduction is \emph{driver-agnostic}: any
analysis that produces a field can be post-processed this way.}
\label{fig:energy-bars}
\end{figure}

\subsection{Eigenfrequency and quality factor, per mode}
\label{sec:fq-table}

The eigenmode driver hands back a family of converged eigenpairs
$(\lambda_m,\vx_m)$ (\cref{ch:eigenvalues}). The \emph{eigenfrequency--$Q$}
reduction turns each into a row $(f_m,Q_m)$ of a per-mode table. Two pieces of
work happen per row, and both \emph{undo} a transform the solve stage applied.

First, the \emph{eigenvalue un-transform}. \Cref{ch:eigenvalues} bent the
spectrum with a shift to make the wanted low eigenvalues easy to find; the
eigenvalue the library returns is therefore in transformed coordinates, and we
must map it back to a physical angular frequency $\omega_m$. For the lossless
\emph{linear} problem $\lambda=\omega^2$, so $\omega=\sqrt\lambda$; for the lossy
\emph{quadratic} problem $\lambda=i\omega$, so $\omega=\lambda/i=-i\lambda$. The
\emph{eigenfrequency} is the real part of the recovered $\omega$:
\begin{equation}
  f_m \;=\; \Re\,\omega_m .
  \label{eq:eigfreq}
\end{equation}

Second, the \emph{quality factor}. $Q$ measures how slowly a mode loses
energy---the ratio of energy stored to energy dissipated per cycle, equivalently
the ratio of the oscillation rate to the decay rate. In the lossy problem the
imaginary part of $\omega$ \emph{is} the decay rate (the loss rate $\kappa_m$),
so
\begin{equation}
  Q_m \;=\; \frac{\omega_m}{\kappa_m},
  \qquad
  \kappa_m = \text{(loss rate from }\Im\,\omega_m\text{)} .
  \label{eq:qfactor}
\end{equation}
A lossless mode has $\kappa_m=0$ and rings forever; its $Q$ is infinite. That is
not a numerical accident to be papered over but a physical fact, and the
reduction guards it explicitly: the verb \code{eigenfreq\_qfactor\_reduce}
returns $Q=\infty$ when $\kappa=0$ rather than dividing by zero.
\begin{lstlisting}[style=l4style]
eigenfreq_qfactor_reduce ptype kappa eigs =
  [ let omega = untransform ptype lambda        -- omega = sqrt(mu) | lambda/i
        f     = re omega                         -- eigenfrequency  f_m = Re omega
        k     = kappa mode                       -- loss rate kappa_m
        q     = if k == 0 then infinity else f / abs k  -- Q = omega/kappa, guarded
    in  (f, q)
  | (lambda, mode) <- eigs ]                      -- map over modes, no inter-mode state
\end{lstlisting}
The two halves are the two scalar projections of each mode; the mode-\emph{field}
recovery ($\vx_m$ itself, and $\vB=-\Curl\vx_m/(i\omega)$) is a separate field
readout, not part of this scalar table. \Cref{fig:omega-plane} shows where $f$
and $Q$ live in the complex $\omega$-plane.

\begin{figure}[t]
\centering
\begin{tikzpicture}[font=\footnotesize, scale=1.0]
  % complex omega-plane
  \draw[->, simGray] (-0.4,0) -- (5.2,0) node[right,font=\scriptsize]{$\Re\,\omega$ (frequency)};
  \draw[->, simGray] (0,-0.4) -- (0,2.7) node[above,font=\scriptsize]{$\Im\,\omega$ (loss rate $\kappa$)};
  % the eigenvalue omega_m
  \coordinate (w) at (3.8,1.1);
  \fill[simRust] (w) circle (2.2pt);
  \node[font=\scriptsize, simRust, anchor=south west] at (w) {$\omega_m$};
  % vector from origin
  \draw[-{Stealth[length=2.2mm]}, thick, simBlue] (0,0) -- (w);
  % drop to real axis: f
  \draw[densely dashed, simGreen, thick] (w) -- (3.8,0);
  \node[font=\scriptsize, simGreen] at (1.9,-0.32) {$f_m=\Re\,\omega_m$};
  % horizontal to imag axis: kappa
  \draw[densely dashed, simRust, thick] (w) -- (0,1.1);
  \node[font=\scriptsize, simRust, anchor=east] at (0,1.1) {$\kappa_m$};
  % distance-to-real-axis annotation
  \node[font=\scriptsize\itshape, simGray, align=left] at (4.2,2.1)
    {$Q_m=\dfrac{\Re\,\omega_m}{\Im\,\omega_m}$:\\high $Q$ $=$\\close to real axis};
  % an arc indicating small angle = high Q
  \draw[simGray, densely dotted] (1.2,0) arc (0:16:1.2);
\end{tikzpicture}
\caption[Eigenfrequency and $Q$ in the complex plane]{Reading $f$ and $Q$ from a
complex eigenvalue. The recovered angular frequency $\omega_m$ is a point in the
complex plane: its \emph{real} part is the oscillation frequency
$f_m=\Re\,\omega_m$ (\cref{eq:eigfreq}), its \emph{imaginary} part is the loss
rate $\kappa_m$ (the decay). The quality factor $Q_m=\Re\,\omega_m/\Im\,\omega_m$
is large precisely when $\omega_m$ sits \emph{close to the real axis}---a small
imaginary part means slow decay, a sharp long-ringing resonance. A purely real
$\omega_m$ (on the axis) is lossless: $\kappa_m=0$, $Q_m=\infty$, the guarded
case in the verb.}
\label{fig:omega-plane}
\end{figure}

\begin{workedexample}{Eigenfrequency and $Q$ from a complex eigenvalue}{fq}
A lossy eigenmode solve returns a converged eigenpair whose recovered angular
frequency is the complex number
\[
  \omega_m \;=\; \omega_r + i\,\omega_i
            \;=\; 2\pi\,(\SI{5.000}{\giga\hertz}) \;+\; i\,(2\pi\cdot\SI{2.5}{\mega\hertz}).
\]
That is, the real part corresponds to a \SI{5}{\giga\hertz} oscillation and the
imaginary part to a \SI{2.5}{\mega\hertz} decay rate. We want the eigenfrequency
$f_m$ and the quality factor $Q_m$.

\emph{Eigenfrequency.} The frequency is the real part, divided by $2\pi$ to go
from angular to ordinary frequency:
\[
  f_m \;=\; \frac{\Re\,\omega_m}{2\pi} \;=\; \SI{5.000}{\giga\hertz}.
\]

\emph{Quality factor.} With the energy/loss convention $Q=\Re\,\omega/(2\,\Im\,\omega)$
(equivalently $\omega_r/(2\omega_i)$, the form in which $Q$ counts radians of
oscillation per e-folding of amplitude),
\[
  Q_m \;=\; \frac{\omega_r}{2\,\omega_i}
        \;=\; \frac{2\pi\cdot 5.000\times 10^{9}}{2\cdot 2\pi\cdot 2.5\times 10^{6}}
        \;=\; \frac{5.000\times 10^{9}}{5.0\times 10^{6}}
        \;=\; 1000 .
\]
A $Q$ of $1000$ is a respectably sharp resonance: the mode rings for about $1000$
radians---some $160$ full cycles---before its amplitude decays by a factor of
$e$. The factor of $2$ in the denominator is the standard convention relating the
amplitude decay rate to the \emph{energy} decay rate (energy goes as amplitude
squared, so it decays twice as fast); different texts place the $2$ differently,
which is why the verb takes the loss-rate convention $\kappa$ as a parameter
rather than hard-coding it.

\emph{The lossless limit.} Now let the loss vanish, $\omega_i\to 0$. The mode
sits exactly on the real axis of \cref{fig:omega-plane}, the denominator goes to
zero, and $Q_m\to\infty$. A lossless cavity rings forever. The reduction does not
divide by zero and emit a \texttt{NaN}; it tests $\kappa=0$ and returns
$Q=\infty$, the honest physical answer. This guard is the one piece of
control flow in an otherwise-arithmetic reduction, and it exists because the
infinite-$Q$ case is real, not exceptional---an ideal lossless resonator is the
limit every low-loss design approaches.
\end{workedexample}

\subsection{Waveguide propagation modes, per mode}
\label{sec:wg-table}

The boundary-mode driver (\cref{ch:waveports}) solves a small two-dimensional
eigenproblem on a port's cross-section and hands back converged eigenpairs. The
\emph{waveguide-mode} reduction turns each into a table row---but this is the one
rank-1 table whose rows carry \emph{fields}, not just scalars. Per mode it
produces:
\begin{itemize}[nosep]
  \item the \emph{propagation constant} $k_n$, recovered from the eigenvalue by
  the same shift-invert un-transform the eigenmode reduction uses
  (\cref{ch:eigenvalues})---it says how fast the mode travels down the guide, and
  whether it travels at all;
  \item the \emph{effective index} $n_{\text{eff}}=k_n/\omega$, the propagation
  constant normalized by the operating frequency (a dimensionless ratio
  comparing the mode's phase velocity to the speed of light);
  \item the \emph{transverse fields} $(\vE_t,\vE_n)$, recovered from the
  eigenvector and \emph{power-normalized} so the mode carries unit Poynting power
  ($\abs{P}=1$); and, for propagating modes only, the longitudinal magnetic field
  \begin{equation}
    \vB_z \;=\; \frac{\Curl\vE_t}{i\,\omega},
    \label{eq:Bz}
  \end{equation}
  the curl of the transverse electric field divided by $i\omega$---a direct
  discrete reading of Faraday's law (\cref{ch:maxwell}, eq.~\eqref{eq:faraday}).
\end{itemize}
The verb is \code{waveguide\_mode\_reduce}; like its scalar-only siblings it is a
\texttt{map} over the converged modes with no inter-mode state, and the operating
frequency $\omega$ rides as a fixed parameter (the $n_{\text{eff}}$ divisor and
the $\vB_z$ scale), not a per-mode datum. The conditional $\vB_z$---present for
propagating modes, absent for evanescent ones---is the one place a table row's
shape varies by content; the framework records it as an optional field.
\Cref{fig:wg-table} sketches such a table, fields and all.

\begin{figure}[t]
\centering
\begin{tikzpicture}[font=\footnotesize]
  % header
  \node[font=\scriptsize\bfseries] at (0,0.55)   {mode};
  \node[font=\scriptsize\bfseries] at (1.4,0.55)  {$k_n$};
  \node[font=\scriptsize\bfseries] at (2.8,0.55)  {$n_{\text{eff}}$};
  \node[font=\scriptsize\bfseries] at (4.9,0.55)  {transverse $\vE_t$};
  \node[font=\scriptsize\bfseries] at (7.0,0.55)  {$\vB_z$};
  \draw[simGray, semithick] (-0.6,0.3) -- (7.7,0.3);
  % --- row 1 (propagating) ---
  \node[font=\scriptsize] at (0,0)   {$1$};
  \node[font=\scriptsize] at (1.4,0) {$31.4$};
  \node[font=\scriptsize] at (2.8,0) {$1.50$};
  \node[font=\scriptsize] at (7.0,0) {\checkmark};
  \draw[simGray, thin] (4.3,-0.25) rectangle (5.5,0.25);
  \draw[-{Stealth[length=1.3mm]}, simRust] (4.55,-0.13) -- (4.55,0.13);
  \draw[-{Stealth[length=1.3mm]}, simRust] (4.85,-0.13) -- (4.85,0.13);
  \draw[-{Stealth[length=1.3mm]}, simRust] (5.15,-0.13) -- (5.15,0.13);
  % --- row 2 (propagating) ---
  \node[font=\scriptsize] at (0,-1.05)   {$2$};
  \node[font=\scriptsize] at (1.4,-1.05) {$22.1$};
  \node[font=\scriptsize] at (2.8,-1.05) {$1.06$};
  \node[font=\scriptsize] at (7.0,-1.05) {\checkmark};
  \draw[simGray, thin] (4.3,-1.30) rectangle (5.5,-0.80);
  \draw[-{Stealth[length=1.3mm]}, simBlue] (4.55,-1.18) -- (4.70,-0.92);
  \draw[-{Stealth[length=1.3mm]}, simBlue] (4.90,-1.18) -- (5.05,-0.92);
  \draw[-{Stealth[length=1.3mm]}, simBlue] (5.20,-1.18) -- (5.30,-0.92);
  % --- row 3 (evanescent) ---
  \node[font=\scriptsize] at (0,-2.10)   {$3$};
  \node[font=\scriptsize] at (1.4,-2.10) {$i\,8.0$};
  \node[font=\scriptsize] at (2.8,-2.10) {---};
  \node[font=\scriptsize] at (7.0,-2.10) {---};
  \node[font=\scriptsize\itshape, simGray] at (4.9,-2.10) {(evanescent: no field)};
  % footnote
  \node[font=\scriptsize\itshape, simGreen, anchor=west] at (-0.6,-2.85)
    {modes 1--2 propagate ($k_n$ real); mode 3 is evanescent ($k_n$ imaginary)};
\end{tikzpicture}
\caption[A waveguide-mode table carrying fields]{The waveguide-mode reduction
produces a rank-1 table whose rows carry \emph{fields}. Each converged mode
yields a propagation constant $k_n$, an effective index
$n_{\text{eff}}=k_n/\omega$, the power-normalized transverse field pattern
$\vE_t$ (the little arrow cartoons), and---for propagating modes only---the
longitudinal field $\vB_z=\Curl\vE_t/(i\omega)$. Modes $1$ and $2$ propagate
($k_n$ real, $n_{\text{eff}}$ defined, $\vB_z$ present, marked \checkmark); mode
$3$ is \emph{evanescent} ($k_n$ imaginary), carries no propagating power, and has
no $\vB_z$. The conditional $\vB_z$ is the one place a table row's shape depends
on its content---recorded as an optional field. This is the only rank-1 table
whose entries are not purely scalar.}
\label{fig:wg-table}
\end{figure}

\section{The taxonomy, and what it buys}
\label{sec:taxonomy}

We have now met all three shapes in full, and the organizing payoff of the
chapter can be stated plainly. \emph{Every} user-facing physical output Palace
produces is one of three reduction shapes (\cref{fig:red-taxonomy}):

\begin{itemize}[nosep]
  \item the \textbf{rank-2 symmetric Gram} $G_{ij}=w(i,j)\,\vx_j^{\!\top}\Kop\vx_i$
  ---capacitance and inductance, one verb differing only in the weight $w$;
  \item the \textbf{rank-2 port-projection}
  $S_{ij}=\sigma_{ij}(\inner{\vs_i}{\vE_j}-\delta_{ij})$ ---S-parameters, a
  linear projection, complex, deliberately not a Gram (\cref{pit:gram-not-proj});
  \item the \textbf{rank-1 scalar-or-mode table}---energy and participation
  (per domain), eigenfrequency and $Q$ (per mode), and waveguide modes (per mode,
  carrying fields), each a pure \texttt{map} with no pairing.
\end{itemize}

\begin{table}[t]
\centering
\caption[The reduction shapes and their products]{The three reduction shapes and
the products each yields. The \emph{shape} column is the structural taxonomy; the
\emph{verb} column is the framework combinator; the \emph{per-item} column is the
computation mapped over the solution family. Reading down the \emph{shape} column,
two reductions produce matrices (Gram, projection) and four produce
tables---and the matrix pair is two genuinely different reductions, not one.}
\label{tab:red-shapes}
\small
\setlength{\tabcolsep}{4pt}
\renewcommand{\arraystretch}{1.25}
\begin{tabular}{@{}llll@{}}
\toprule
Shape & Verb & Per-item computation & Product \\
\midrule
rank-2 Gram & \code{gram\_reduce} & $w(i,j)\,\vx_j^{\!\top}\Kop\vx_i$ & capacitance \\
rank-2 Gram & \code{gram\_reduce} & $w{=}1/(I_iI_j)$, same form & inductance \\
rank-2 projection & \code{sparameter\_reduce} & $\sigma_{ij}(\inner{\vs_i}{\vE_j}-\delta_{ij})$ & S-parameters \\
rank-1 table & \code{domain\_energy\_reduce} & $(\tfrac12\inner{f}{\mat{M}_if},\,p_i)$ & energy, participation \\
rank-1 table & \code{eigenfreq\_qfactor\_reduce} & $(\Re\,\omega_m,\,\omega_m/\kappa_m)$ & eigenfrequency, $Q$ \\
rank-1 table & \code{waveguide\_mode\_reduce} & $(k_n,n_{\text{eff}},\vE_t,\vB_z)$ & waveguide modes \\
\bottomrule
\end{tabular}
\end{table}

\Cref{tab:red-shapes} collects them. The lesson is the book's recurring refrain,
now at the output stage: a small handful of shapes, specialized by a weight, a
basis, or a per-item formula, produce the entire menu of physical answers. The
taxonomy is not a tidy-minded over-unification---we were careful in
\cref{pit:gram-not-proj} to keep the Gram and the projection distinct---but a
genuine structural economy: there really are only three ways the machinery
collapses a field family into an answer.

\begin{keyidea}
The \textbf{three reduction shapes} produce \emph{every} physical output:
the rank-2 symmetric Gram (capacitance, inductance), the rank-2 port-projection
(S-parameters), and the rank-1 scalar-or-mode table (energy/participation,
eigenfrequency/$Q$, waveguide modes). Two shapes yield matrices and one yields
tables; the two matrix reductions are \emph{different} reductions that merely
share an output rank. Each is an embarrassingly parallel \texttt{map} over the
solution family---the last and simplest stage of every analysis, and the place
where the physics re-enters as a capacitance, a scattering coefficient, or a
resonant frequency.
\end{keyidea}

\section{Closing Part IV: the machinery is assembled}
\label{sec:bridge}

This chapter closes Part~IV, and with it the long climb up the machinery. It is
worth standing back to see how far we have come, and where Part~V takes us.

Part~II built the \emph{field theory}: Maxwell's equations, the weak form, the
finite-element function spaces, and the assembly that turns continuous physics
into discrete operators $\Kop,\Mop,\Cop$. Part~III built the \emph{calculus}: the
abstract vocabulary of tensors, records, folds, combinators, monads, and
operators in which the whole simulator can be written down precisely. Part~IV
built the \emph{numerical machinery}: the Krylov solvers (\cref{ch:krylov,ch:cg,ch:gmres}),
their orthogonalization and rotation kernels (\cref{ch:orthog,ch:givens}), the
preconditioners and multigrid that make them converge (\cref{ch:precond,ch:smoothers,ch:multigrid}),
the eigenvalue machinery and its spectral transforms (\cref{ch:eigenvalues,ch:krylovschur,ch:deflation}),
the time integrators (\cref{ch:time}), the adaptive-refinement outer loop
(\cref{ch:amr}), and now the output reductions that turn solved fields into
physical products.

Every box of the skeleton (\cref{ch:situating}) is now filled with real
machinery. What remains is to \emph{compose} it. That is Part~V---the
\emph{modalities}. Each modality chapter takes one of the six drivers and walks
its full path through the assembled machinery: which function space it assembles
on, which solve corner it occupies, which solver it calls, and---the subject of
this chapter---which reduction shape produces its answer. The electrostatic
driver (\cref{ch:electrostatics}) composes a fixed-operator family solve with a
unit-weight Gram; the magnetostatic (\cref{ch:magnetostatics}) swaps the space
and the Gram weight; the driven (\cref{ch:driven}) swaps in an operator-varying
sweep and the port-projection; the eigenmode (\cref{ch:eigenmodes}) calls the
black-box eigensolver and reduces to an $(f,Q)$ table; the boundary-mode
(\cref{ch:waveports}) does the same on a cross-section and reduces to a waveguide
table. The reductions of this chapter are the last stage of each of those paths.
The parts are all on the bench; Part~V assembles the machines.

\summarysection
The \texttt{reduce} stage turns a driver's solution \emph{family} into the
user-facing physical product. Every reduction is an embarrassingly parallel
\texttt{map} over the family with no inter-element state---never a fold
(\cref{ch:combinators}). There are exactly three reduction \emph{shapes}. The
\textbf{rank-2 symmetric Gram} $G_{ij}=w(i,j)\,\vx_j^{\!\top}\Kop\vx_i$
(\code{gram\_reduce}) produces capacitance ($w=1$, per-terminal potentials) and
inductance ($w=1/(I_iI_j)$, per-source vector potentials) from one verb, its
symmetry structural and its diagonal twice the stored field energy. The
\textbf{rank-2 port-projection}
$S_{ij}=\sigma_{ij}(\inner{\vs_i}{\vE_j}-\delta_{ij})$
(\code{sparameter\_reduce}) produces S-parameters by projecting each driven field
onto the port modes; it is complex, non-symmetric, and deliberately \emph{not} a
Gram---sharing an output rank is not sharing a meaning. The \textbf{rank-1
scalar-or-mode table} produces energy and participation per domain
(\code{domain\_energy\_reduce}, driver-agnostic), eigenfrequency and $Q$ per mode
(\code{eigenfreq\_qfactor\_reduce}, with the $\kappa{=}0\Rightarrow Q{=}\infty$
guard), and waveguide modes per mode (\code{waveguide\_mode\_reduce}, the one
table carrying fields). These three shapes, specialized by a weight, a basis, or
a per-item formula, produce the simulator's entire menu of physical answers---and
they are the last stage of each modality in Part~V.

\furtherreading
The physical meanings of the reduced quantities are standard electromagnetics:
Griffiths~\citep{griffiths2017} for capacitance, inductance, and stored field
energy (the Gram diagonal); Pozar~\citep{pozar2011} for S-parameters, quality
factors, and waveguide propagation modes (the microwave quantities, including the
de-embedding and impedance scalings of the port projection). For the
finite-element side---how the energy operators $\Kop,\mat{M}_i$ are assembled and
why the diagonal of the Gram is exactly twice the discrete field energy---Jin's
finite-element electromagnetics text~\citep{jin2014} is the natural companion,
and it treats the divergence-free projection and the spurious-mode cleanup that
keep the eigenfrequency table honest. The structural observation that all of
these are one of three \texttt{map}-shaped reductions is this book's own, distilled
from reading Palace's postprocessing side by side; \cref{ch:situating} states the
taxonomy and this chapter fills it in.

\exercisessection
\begin{enumerate}[nosep]
  \item \textbf{One verb, two products.} Write down the Gram reduction entry
  $G_{ij}$ and the weight $w(i,j)$ for capacitance and for inductance. In one
  sentence each, explain why capacitance's weight is trivial ($w=1$) but
  inductance's ($w=1/(I_iI_j)$) is not.
  \item \textbf{Mirror, don't recompute.} A three-terminal capacitance problem
  needs the full $3\times 3$ matrix $\mat{C}$. How many distinct bilinear-form
  evaluations $\vx_j^{\!\top}\Kop\vx_i$ does \code{gram\_reduce} actually perform,
  and which entries are filled by \code{symmetric\_from\_upper} instead? What
  property of $\Kop$ licenses the mirroring?
  \item \textbf{Sign of the off-diagonal.} In \cref{wex:cap2x2} the mutual
  capacitance $C_{12}$ came out negative. Explain physically why the off-diagonal
  of a Maxwell capacitance matrix is always negative, using the relation
  $Q_1=C_{11}V_1+C_{12}V_2$.
  \item \textbf{Not a Gram.} Give the three structural reasons (from
  \cref{pit:gram-not-proj}) that the S-parameter matrix is not a Gram matrix,
  even though both are square and port-indexed. For each reason, name the specific
  feature of \code{sparameter\_reduce}'s body that a unified ``matrix reduction''
  verb would have to violate.
  \item \textbf{The $Q$ guard.} Recompute $Q_m$ from \cref{wex:fq} if the loss
  rate were $\omega_i=2\pi\cdot\SI{5}{\mega\hertz}$ instead of \SI{2.5}{\mega\hertz}.
  Then take $\omega_i\to 0$: what does the reduction return, and why is that the
  physically correct answer rather than a numerical failure?
  \item \textbf{Two undone transforms.} The eigenfrequency--$Q$ reduction
  ``un-transforms'' the eigenvalue twice over. Name the two transforms it undoes:
  one applied by the eigenvalue \emph{solve} stage (\cref{ch:eigenvalues}) and one
  implicit in the lossy problem's eigenvalue convention $\lambda=i\omega$. What
  does each contribute to the final $(f,Q)$ row?
  \item \textbf{Driver-agnostic energy.} The domain-energy reduction is the one
  reduction that runs identically as a post-process of \emph{any} field-bearing
  analysis. State the minimal precondition it needs from the driver above it, and
  explain why the participation factors sum to one only when the configured
  domains partition the field's support.
  \item \textbf{Place the products.} For each of the six Palace analyses, name its
  product and the reduction shape (Gram / projection / table) that produces it.
  Which shape is used by the most analyses, and why is that unsurprising given
  that a table is just a \texttt{map} with no pairing?
\end{enumerate}
